\documentclass{beamer}
\usepackage{fontspec}
\usepackage{bidi}
\usepackage{polyglossia}

\setdefaultlanguage{hebrew}
\setotherlanguage{english}
\newfontfamily\hebrewfont{Arial}

\usetheme{Madrid}

\begin{document}

\begin{frame}
  \begin{center}
    {\Large \textbf{פיתוח מערכת מבוססת \LRE{FPGA} להתראת נפילה לחולי דמנציה}} \\ \vspace{0.5cm}
    קליטת תמונה בזמן אמת ממצלמת \LRE{OV7670} והצגתה דרך כרטיס \LRE{Artix-7}. \\
    על מנת לאפשר עיבוד תמונה לזיהויי עמידה ממושכת לחולי דמנציה. \\ \vspace{0.5cm}
    \textit{מגיש: יוסי ברים | הנדסת חשמל שנה ד | המרכז האקדמי לב}
  \end{center}
\end{frame}

\begin{frame}{ארכיטקטורת המערכת}
  \begin{itemize}
    \item נתוני התמונה הגולמיים נקלטים מהמצלמה ונאגרים בזיכרון ה-\LRE{BRAM}.
    \item בקר ה-\LRE{VGA} שולף את הנתונים ומעביר לממיר \LRE{D to A} לקבלת אות אנלוגי פיזי למסך.
  \end{itemize}
\end{frame}

\begin{frame}{תזמוני VGA: הפיזיקה מאחורי המספרים}
  \begin{table}[]
    \begin{tabular}{|l|c|r|r|}
    \hline
    \textbf{פרמטר} & \textbf{שעונים} & \textbf{משמעות פיזיקלית} & \textbf{מימוש (RTL)} \\ \hline
    \LRE{Active Display} & 640 & הקרן מציירת פיקסלים. & שידור \LRE{RGB} פעיל. \\ \hline
    \LRE{Front Porch} & 16 & מרווח לקרן להירגע. & כיבוי ל-\LRE{0000}. \\ \hline
    \end{tabular}
  \end{table}
\end{frame}

\end{document}
